\uuid{N94v}
\titre{Matrices, déterminants et inversion}
\niveau{L1}
\module{Algèbre}
\chapitre{Matrice}
\sousChapitre{Inverse, méthode de Gauss-Jordan}
\theme{Calcul de déterminants, matrice de codage, application linéaire}
\auteur{}
\datecreate{2026-06-04}
\organisation{AMSCC}
\difficulte{2}

\contenu{
	\texte{
		Un système de communication encode des messages en transformant des vecteurs
		$V \in \mathbb{R}^3$ en vecteurs codés $V' = AV$, où :
		\[
		A = \begin{pmatrix}1&0&1\\0&2&1\\1&1&2\end{pmatrix}.
		\]
		Décoder un message consiste donc à retrouver $V$ à partir de $V'$.
	}
	\begin{enumerate}
		
		\item
		\question{Calculer $\det(A)$. Que peut-on en conclure sur l'inversibilité de $A$ et sur la possibilité de décoder les messages ?}
		\reponse{
			On développe par rapport à la première ligne :
			\[
			\det(A)
			=
			1\begin{vmatrix}2&1\\1&2\end{vmatrix}
			-
			0\begin{vmatrix}0&1\\1&2\end{vmatrix}
			+
			1\begin{vmatrix}0&2\\1&1\end{vmatrix}.
			\]
			Ainsi :
			\[
			\det(A)
			=
			(4-1) + (0-2)
			=
			3-2
			=
			1.
			\]
			Donc $\det(A)=1\neq 0$, ce qui montre que $A$ est inversible.
			
			Le système de codage est donc décodable : pour tout vecteur codé $V'$, il existe un unique vecteur initial $V$ tel que $V'=AV$.
		}
		
		\item
		\question{Calculer la matrice inverse $A^{-1}$ à l'aide de la méthode de Gauss-Jordan.}
		\reponse{
			On applique la méthode de Gauss-Jordan à la matrice augmentée $\left(A \mid I_3\right)$ :
			\[
			\left(
			\begin{array}{ccc|ccc}
				1&0&1&1&0&0\\
				0&2&1&0&1&0\\
				1&1&2&0&0&1
			\end{array}
			\right).
			\]
			
			On effectue d'abord l'opération $L_3 \leftarrow L_3-L_1$ :
			\[
			\left(
			\begin{array}{ccc|ccc}
				1&0&1&1&0&0\\
				0&2&1&0&1&0\\
				0&1&1&-1&0&1
			\end{array}
			\right).
			\]
			
			On échange $L_2$ et $L_3$ :
			\[
			\left(
			\begin{array}{ccc|ccc}
				1&0&1&1&0&0\\
				0&1&1&-1&0&1\\
				0&2&1&0&1&0
			\end{array}
			\right).
			\]
			
			On effectue $L_3 \leftarrow L_3-2L_2$ :
			\[
			\left(
			\begin{array}{ccc|ccc}
				1&0&1&1&0&0\\
				0&1&1&-1&0&1\\
				0&0&-1&2&1&-2
			\end{array}
			\right).
			\]
			
			On remplace $L_3$ par $-L_3$ :
			\[
			\left(
			\begin{array}{ccc|ccc}
				1&0&1&1&0&0\\
				0&1&1&-1&0&1\\
				0&0&1&-2&-1&2
			\end{array}
			\right).
			\]
			
			On annule ensuite les coefficients au-dessus du pivot de la troisième colonne :
			$L_1 \leftarrow L_1-L_3$, $L_2 \leftarrow L_2-L_3$. On obtient :
			\[
			\left(
			\begin{array}{ccc|ccc}
				1&0&0&3&1&-2\\
				0&1&0&1&1&-1\\
				0&0&1&-2&-1&2
			\end{array}
			\right).
			\]
			
			Ainsi :
			\[
			A^{-1}
			=
			\begin{pmatrix}
				3&1&-2\\
				1&1&-1\\
				-2&-1&2
			\end{pmatrix}.
			\]
		}
		
		\item
		\question{Décoder les trois vecteurs suivants en calculant les vecteurs originaux
			$V_1$, $V_2$, $V_3$ :
			\[
			V'_1 = \begin{pmatrix}4\\7\\9\end{pmatrix},\quad
			V'_2 = \begin{pmatrix}3\\3\\5\end{pmatrix},\quad
			V'_3 = \begin{pmatrix}4\\1\\5\end{pmatrix}.
			\]}
		\reponse{
			Puisque $V'=AV$, on obtient $V=A^{-1}V'$. On calcule donc :
			\begin{align*}
				V_1
				&=
				\begin{pmatrix}
					3&1&-2\\
					1&1&-1\\
					-2&-1&2
				\end{pmatrix}
				\begin{pmatrix}4\\7\\9\end{pmatrix}
				=
				\begin{pmatrix}
					12+7-18\\
					4+7-9\\
					-8-7+18
				\end{pmatrix}
				=
				\begin{pmatrix}1\\2\\3\end{pmatrix},
				\\[6pt]
				V_2
				&=
				\begin{pmatrix}
					3&1&-2\\
					1&1&-1\\
					-2&-1&2
				\end{pmatrix}
				\begin{pmatrix}3\\3\\5\end{pmatrix}
				=
				\begin{pmatrix}
					9+3-10\\
					3+3-5\\
					-6-3+10
				\end{pmatrix}
				=
				\begin{pmatrix}2\\1\\1\end{pmatrix},
				\\[6pt]
				V_3
				&=
				\begin{pmatrix}
					3&1&-2\\
					1&1&-1\\
					-2&-1&2
				\end{pmatrix}
				\begin{pmatrix}4\\1\\5\end{pmatrix}
				=
				\begin{pmatrix}
					12+1-10\\
					4+1-5\\
					-8-1+10
				\end{pmatrix}
				=
				\begin{pmatrix}3\\0\\1\end{pmatrix}.
			\end{align*}
		}
		
		\item
		\question{On remplace maintenant la matrice de codage par
			\[
			B = \begin{pmatrix}1&2&1\\2&4&2\\1&3&2\end{pmatrix}.
			\]
			Expliquer pourquoi ce nouveau système de codage serait défaillant.}
		\reponse{
			On remarque que la deuxième ligne de $B$ est égale au double de la première :
			$(2,4,2)=2(1,2,1)$.
			Les lignes de $B$ sont donc linéairement dépendantes, par conséquent $\det(B)=0$
			et $B$ n'est pas inversible.
			
			Le décodage n'est alors pas possible de manière unique : il n'existe pas de matrice
			$B^{-1}$ permettant de retrouver systématiquement $V$ à partir de $V'$.
			Deux messages distincts peuvent produire le même vecteur codé ; le système est donc
			défaillant.
		}
		
	\end{enumerate}
}